% ===============================================================
% CHAPITRE 1 : CONTEXTE ET PROBLÉMATIQUE
% ===============================================================

\chapter{Contexte et Problématique}

\section*{Introduction}

Ce chapitre présente le contexte général dans lequel s'inscrit le projet WC Fantasy, définit la problématique à résoudre et précise les objectifs poursuivis. Il constitue le fondement sur lequel repose l'ensemble de la démarche de conception et de développement.

\section{Contexte du projet}

Le football est le sport le plus populaire au monde avec \textbf{plus de 5 milliards de fans} selon la FIFA \cite{fifa_wc2026}. La Coupe du Monde 2026 représente un événement sans précédent : pour la première fois dans l'histoire de la compétition, \textbf{48 nations s'affronteront}, générant \textbf{104 matchs} sur un mois et demi. Cet événement crée une opportunité unique pour des applications de divertissement sportif.

Le fantasy sport est un secteur en plein essor qui permet aux supporters de s'impliquer activement dans leur passion. Le principe est simple : chaque utilisateur constitue une \textbf{équipe virtuelle} composée de joueurs réels, et les performances réelles de ces joueurs (buts, passes décisives, clean sheets, etc.) se traduisent en points dans le jeu \cite{statista_fantasy}. Ces points déterminent un \textbf{classement} au sein de ligues privées ou publiques.

En Afrique du Nord et au Maroc en particulier, l'intérêt pour le football est exacerbé par la \textbf{première participation historique du Maroc à la co-organisation} d'une Coupe du Monde, créant un terreau extrêmement favorable pour ce type d'application. L'essor de la téléphonie mobile (\textbf{taux de pénétration supérieur à 80\% au Maroc}) renforce cette opportunité.

\section{Problématique}

Malgré la popularité du fantasy sport, plusieurs problèmes persistent dans les solutions actuellement disponibles :

\begin{itemize}
    \item \textbf{Absence d'application dédiée} à la Coupe du Monde 2026 avec un catalogue complet de joueurs des 48 nations participantes.
    \item Les plateformes existantes sont \textbf{peu adaptées} aux compétitions éphémères (durée limitée, format à élimination directe).
    \item \textbf{Manque de fonctionnalités sociales} : chat en temps réel entre membres d'une ligue, batailles directes (head-to-head), prédictions communautaires.
    \item \textbf{Interfaces mobiles souvent non optimisées}, lentes et peu gamifiées.
    \item \textbf{Absence d'intelligence artificielle} pour enrichir l'expérience utilisateur (recommandations, chatbot d'assistance).
\end{itemize}

La question centrale est donc : \textbf{\textit{comment concevoir une application mobile de fantasy football complète, performante et engageante, spécifiquement adaptée au format et à l'échelle de la Coupe du Monde 2026 ?}}

\section{Objectifs du projet}

\textbf{Objectif général :} Concevoir et développer une application mobile de fantasy football dédiée à la Coupe du Monde 2026, offrant une expérience utilisateur complète, sécurisée et en temps réel.

\subsection*{Objectifs spécifiques :}

\begin{enumerate}
    \item \textbf{Offrir une expérience utilisateur immersive} avec une interface moderne, réactive et facile d'utilisation pour tous les profils de supporters.
    \item \textbf{Créer un environnement hautement social} permettant aux utilisateurs de créer des ligues privées, de communiquer en temps réel et de s'affronter directement (Head-to-Head).
    \item \textbf{Assurer un suivi en temps réel} des performances des joueurs lors de la compétition, avec un calcul de points instantané.
    \item \textbf{S'adapter au format court d'une Coupe du Monde} avec une dynamique de jeu nerveuse, des prédictions de matchs et des récompenses régulières.
\end{enumerate}

\section*{Conclusion}

Ce chapitre a permis de situer le projet WC Fantasy dans son contexte technique, social et économique. La problématique clairement identifiée — absence d'une plateforme mobile dédiée à la Coupe du Monde 2026 avec des fonctionnalités sociales riches — justifie pleinement la conception d'une solution sur mesure. Le chapitre suivant présente une étude commerciale du marché et analyse les solutions existantes (leurs forces, faiblesses et les retours des utilisateurs) afin de définir le positionnement de notre projet.
